\chapter{Concluzii}

\section{Rezultate ob\c tinute}

Lucrarea de licen\c t\u a a avut ca obiectiv dezvoltarea \textbf{aplica\c tiei web pentru organizarea \c si gestionarea evenimentelor}, orientate \^int\^ai spre conferin\c te, meeting-uri, concerte \c si evenimente culturale, cu suport complementar pentru competi\c tii sportive. Obiectivele propuse au fost atinse:

\begin{itemize}
	\item Aplica\c tie full-stack func\c tional\u a: React + Express + PostgreSQL.
	\item Autentificare complet\u a: \^inregistrare cu verificare email, login JWT, resetare parol\u a, profil, \c stergere cont.
	\item Listare \c si filtrare evenimente (nume, jude\c t, dat\u a).
	\item Pl\u a\c ti Stripe Checkout \^in RON; bilete nominale \c si ,,Biletele mele``.
	\item Panou admin: CRUD evenimente cu imagini, check-in, gestionare admini.
	\item Documentare: arhitectur\u a, ERD, UML, API REST, scenarii pe tipuri de evenimente.
\end{itemize}

Solu\c tia demonstreaz\u a integrarea unor tehnologii moderne \^intr-un produs coerent, adaptat contextului rom\^anesc (jude\c te, RON, email Brevo).


\section{Limit\u ari}

\begin{itemize}
	\item Lips\u a teste automate (unit/integration) \c si pipeline CI/CD.
	\item Promovarea primului administrator necesit\u a script SQL sau \texttt{promote-admin}.
	\item Emailurile depind de configurarea SMTP; \^in development, codurile pot ap\u area \^in consola backend.
	\item Confirmarea pl\u a\c tii se bazeaz\u a pe apel frontend \texttt{confirm-payment}; webhook Stripe ar \^imbun\u at\u a\c ti robuste\c tea.
	\item Un singur tip de pre\c t per eveniment — f\u ar\u a categorii VIP/standard sau locuri numerotate.
\end{itemize}

\section{Dezvolt\u ari viitoare}

Dezvoltarea aplica\c tiei poate fi continuat\u a prin abordarea limit\u arilor identificate \c si ad\u augarea de noi func\c tionalit\u a\c ti:

\begin{itemize}
	\item Categorii de bilete (conferin\c t\u a: early bird / regular; concert: VIP).
	\item Export participan\c ti CSV din panoul admin.
	\item Notific\u ari email la cump\u arare bilet reu\c sit\u a.
	\item Paginare evenimente \c si caching.
	\item Webhook Stripe pentru confirmare server-side.
	\item Teste automate \c si deploy containerizat (Docker).
	\item Aplica\c tie mobil\u a sau PWA pentru check-in offline.
	\item Module op\c tionale pentru evenimente sportive (echipe, program).
\end{itemize}

Aplica\c tia reprezint\u a o baz\u a solid\u a pentru extindere \^intr-un produs destinat organizatorilor de conferin\c te, concerte \c si meeting-uri, cu accent pe simplitate, securitate \c si experien\c ta utilizatorului.